% !TEX root = ../main.tex
% Metodologia. Nenhum numero digitado: todo valor vem de uma macro \F... de fatos.tex,
% gerada por artigo/tools/folha_de_fatos.py. O conferidor (folha_de_fatos.py --conferir)
% recusa digito solto neste arquivo.

\section{Study Design}
\label{sec:design}

\subsection{Research Questions}
\label{sec:rqs}

We ask whether the production classes that professional developers perceive as smelly are
tested by worse tests. We operationalize ``worse tests'' as test smells detected by a
static tool, and ``perceived as smelly'' as the severity that reviewers assigned to the
class in a curated code smell data set.

\begin{description}
  \item[RQ1] Is the perceived severity of code smells in a production class associated
  with the test smell density of the tests dedicated to that class?
  \item[RQ2] Is any individual test smell, or any pair of code smell and test smell,
  associated with perceived severity?
  \item[RQ3] How do the choices of outcome normalization and of test-to-code binding
  change the answer to RQ1?
\end{description}

RQ1 is the main question and is answered with three operationalizations
(Section~\ref{sec:analysis}). RQ2 checks whether a null aggregate result hides opposite
effects of individual smells. RQ3 makes explicit the two measurement decisions that, in
our data, decide whether an association appears at all.

Figure~\ref{fig:pipeline} summarizes the pipeline: we collect test smells in the
repositories of the code smell data set, at the commit the reviewers saw
(Section~\ref{sec:collection}); we link each annotated production class to its tests
(Section~\ref{sec:binding}); and we correlate the two sides
(Section~\ref{sec:analysis}).

\begin{figure}[t]
  \centering
  \fbox{\parbox{0.9\linewidth}{\centering\vspace{2em}pipeline: MLCQ $\rightarrow$ mirrors
  at the MLCQ commit $\rightarrow$ JNose $\rightarrow$ binding $\rightarrow$ analysis
  \vspace{2em}}}
  \caption{Study pipeline. Every number reported in this paper is produced by a versioned
  script of the replication package.}
  \label{fig:pipeline}
\end{figure}

\subsection{Code Smell Labels: the MLCQ Data Set}
\label{sec:mlcq}

MLCQ~\cite{madeyski2020mlcq} contains \FrevisoesMlcq{} reviews by \FrevisoresMlcq{}
professional developers of \FamostrasMlcq{} code samples from open-source Java projects.
Each review grades one sample for two smells at once: blob and data class for classes,
feature envy and long method for methods. The grade is a severity on the scale none,
minor, major, critical, which we code from zero to \FescalaMax{}. Our unit of label is
therefore the pair (sample, smell); \FamostrasUmSmell{} samples were graded for only one
of their two smells, so MLCQ has \FunidadesMlcq{} such units.

\paragraph{The sampling design makes binary labels circular}
MLCQ's appendix states that the reviewers of the last phase cross-checked only samples
already graded above none. We verified this in the data: all \FamostrasCrosscheck{}
cross-checked samples had a previous positive grade. The number of reviewers of a unit is
thus a consequence of its label, not independent evidence for it: \FpctUnidadesUmRevisor{}
of the units have a single reviewer, and almost all of them are negative. A rule such as
``positive if any reviewer graded above none'' discards the verdict of the cross-check,
which was collected precisely to judge the flagged samples: \FpctDisputadas{} of the
positive units with two or more reviewers were graded none by some reviewer.

The same design lets us measure how reliable an isolated ``none'' is. Because samples, and
not units, were routed to the cross-check, \FnoneRevisados{} units graded none were
reviewed again only because the other smell of the sample had been flagged. In
\FpctNoneSeSustenta{} of them the isolated ``none'' agreed with the final majority.

\paragraph{Label used}
We avoid any threshold and use the mean severity of the reviewers of each unit as a
continuous label. A production file may hold several samples of the same smell (e.g.,
several methods judged for long method); the file's severity for that smell is the
largest among them. The aggregate label of the file, used in RQ1, is the largest severity
among the smells it was judged for, i.e., the smell perceived as most severe. Smells for
which a file was not judged remain missing and are never treated as zero.

For the group comparisons we also classify each unit by strength of evidence:
\emph{confirmed} (the majority graded above none), \emph{clean} (all reviewers graded
none), and \emph{disputed} (someone graded above none, but not the majority).

\subsection{Test Smell Collection}
\label{sec:collection}

\paragraph{Repositories}
MLCQ references \FreposMlcq{} repositories, each at a specific commit. We use mirrors of
these repositories created at that commit and match each one by commit hash. We analyze
the \FreposHashOk{} repositories whose mirror head equals the MLCQ commit; the collection
script re-checks the head of every clone. The remaining repositories have a mirror at a
different commit (\FreposHashDiverge{}) or no mirror (\FreposSemEspelho{}), and are
excluded: analyzing another commit would pair the reviewers' judgment with code they did
not see.

\paragraph{Detector}
We detect test smells with JNose Test~\cite{virginio2020jnose}, version 2.5.0
(jnose-core 0.9.4), a JUnit-aware extension of tsDetect~\cite{peruma2020tsdetect}. We use
the tool's default configuration of \FsmellsHabilitados{} test smells rather than
selecting smells ourselves, with one exception: Unknown Test. Its detector, when it fails
on a class, discards the analysis of all other smells of that class; in a separate run
with the full default configuration it would have erased \FunknownOcorrencias{}
occurrences of other smells in \FunknownClasses{} classes.

The collection covers \FclassesTeste{} test classes in \FreposComTeste{} repositories and
yields \Focorrencias{} test smell occurrences. A file whose name matches JNose's test
pattern and imports JUnit, but declares no top-level class or interface, aborts the
analysis of the whole project without any error on screen. Before analysis, the
collection script removes exactly these files from each clone, as decided by an oracle
that replicates JNose's own checks on the syntax tree with the parser bundled in the tool.
Such a file can never yield a test class, so removing it changes nothing in the
collection; \FarquivosRemovidosNpe{} files in \FreposRemovidosNpe{} repositories were
removed.

\paragraph{A known defect of the detector}
JNose names a test class after the \emph{last} class declaration it visits. In files with
a nested class it reports the nested helper's name, and it resolves the production class
from that name. The two smells that compare the test against its production class, Eager
Test and Lazy Test, are therefore mostly missed: their rate drops \FquedaEager{} and
\FquedaLazy{} in these files, while General Fixture, which reads only the test, is stable.
This affects \FclassesAninhadas{} test
classes (\FpctClassesAninhadas{}). We keep these classes and flag them; RQ2 and the
sensitivity analyses report results with and without them.

\subsection{Test-to-Code Binding}
\label{sec:binding}

To answer RQ1 we need, for each annotated production file, the tests dedicated to it. We
link production files and test classes with three strategies, each recorded per pair so
that its precision can be measured separately:

\begin{itemize}
  \item \emph{Exact path}: the production file that JNose itself resolves for the test
  class. We disassembled the tool to check what this is: a naming convention (strip the
  \texttt{Test} prefix or suffix) followed by a file name search in the project.
  \item \emph{Convention}: \texttt{X.java} is linked to \texttt{XTest},
  \texttt{XTests}, \texttt{XTestCase}, \texttt{XIT}, or \texttt{XITCase} in the same
  repository; a pair is marked as \emph{mirrored} when test and production share the
  package path.
  \item \emph{Static reference}: the test class references the production type, through
  an exact import, a wildcard import, or the same package.
\end{itemize}

The first two strategies form the \emph{deterministic} branch: \FarquivosDet{} production
files in \FreposDet{} repositories, linked to \FclassesTesteDet{} test classes. Adding
static references forms the \emph{extended} branch: \FarquivosAmp{} files in
\FreposAmp{} repositories, linked to \FclassesTesteAmp{} test classes. Production files
referenced by more than \FmaxTestes{} test classes are excluded from this branch, to keep
widely shared types, such as data objects used across a test suite, from dominating it.

\paragraph{Precision audit} We audited \FparesAuditados{} pairs, sampled at random and
stratified by strategy and evidence (\FestratosAuditoria{} strata, up to
\FparesPorEstratoAuditoria{} pairs each; the two convention strata are a census). The
first author judged all pairs, blind to the stratum, under two criteria: whether the
production class is the \emph{target} of the test (strict), and whether the test
\emph{executes} code of the class, even as a fixture (broad). Each pair was judged from an
extract of the two files at the collection commit: the test lines that reference the
class, with their full statements and the calls made on the resulting objects. A pair the
extract could not settle was marked undecided. A second rater, not involved in the study,
independently judged a stratified subsample of \FparesSegundoAvaliador{} pairs
(\FparesPorEstratoSegundo{} per stratum). Inter-rater agreement was
\FconcordanciaEstrito{} (Cohen's $\kappa$ = \FkappaEstrito{}) under the strict criterion
and \FconcordanciaAmplo{} ($\kappa$ = \FkappaAmplo{}) under the broad one, computed before
disagreements were resolved by discussion. The audit protocol was registered before the
first pair was judged.

The precision of each branch is the per-stratum precision weighted by the stratum's share
of the pairs, over the decided pairs. The deterministic branch has a precision of
\FprecisaoHumanaDet{}. All of its \FerrosHumanosDet{} errors are the same case: a test
named after the class that never references it. The extended branch has a precision of
\FprecisaoHumanaAmp{} under the strict criterion and \FprecisaoHumanaAmpAmplo{} under the
broad one (\FprecisaoHumanaAmpAmploInf{} to \FprecisaoHumanaAmpAmploSup{} if every
undecided pair is counted as wrong or as right): a static reference identifies the classes
a test \emph{exercises}, not the class it \emph{tests}. We therefore treat the two
branches as answering different questions. The deterministic branch answers RQ1; the
extended branch is reported for RQ3, as the answer one would get by conflating the two
constructs.

\subsection{Outcome Variables}
\label{sec:outcomes}

For each production file we sum, over its linked test classes, the occurrences of each
test smell. RQ1 uses three outcomes:

\begin{itemize}
  \item \emph{raw}: the total number of test smell occurrences;
  \item \emph{density}: the raw count divided by the number of members of the linked test
  classes;
  \item \emph{distinct}: how many of the \FsmellsHabilitados{} test smells occur at least
  once.
\end{itemize}

Raw and distinct grow with the size of the test. Density is our primary outcome because
it isolates the quality of the test from its size: its correlation with test size is
\FrhoDensidadeTamanho{}, against \FrhoBrutoTamanho{} for the raw count.

The denominator of density is the member count that JNose reports for the class
(methods, fields, constructors, and nested types), not the number of test methods. We
report density as occurrences per class member and do not claim it is per test method.
In classes with a nested type the count is the helper's, which underestimates the
denominator; the sensitivity analysis without these files addresses this.

\subsection{Analysis}
\label{sec:analysis}

The variables are far from normal: most files have zero occurrences of most smells. We
therefore use rank-based statistics throughout.

\paragraph{RQ1}
Spearman's $\rho$ between each outcome and the aggregate label, with 95\% confidence
intervals from Fisher's $z$ transformation. To quantify the role of test size, we also
report the partial Spearman correlation controlling for the member count of the tests,
and the correlation of the label itself with test size. As a second operationalization,
we compare files whose label was confirmed by the reviewers' majority
(\FnConfirmados{} files) with clean files (\FnLimpos{}) using the Mann--Whitney $U$ test,
with the rank-biserial correlation as effect size. As a third, we compare confirmed files
with files that were flagged but not confirmed (\FnDisputados{}). This last comparison is
the only one immune to the sampling design of MLCQ, since both groups were flagged and
routed to the cross-check.

\paragraph{RQ2}
Spearman's $\rho$ for each test smell, in raw and density form, against the aggregate
label; and for each pair of code smell and test smell, in density form, against the
severity of that code smell. Default Test does not occur in any analyzed file and is
excluded as a constant. Each set of tests is a family, adjusted with
Benjamini--Hochberg~\cite{benjamini1995fdr}.

\paragraph{Power and equivalence}
A null result is only informative if the design could have detected an effect. With the
\FarquivosDet{} files of the deterministic branch, a two-sided test at $\alpha = 0.05$ has 80\% power to detect
$|\rho| \geq$ \FrDetectavelDet{}. Beyond power, we report the largest effect compatible
with the data, i.e., the bound of the 90\% confidence interval farthest from zero, and
run two one-sided tests (TOST)~\cite{lakens2017equivalence} against
$|\rho| = $ \FmargemEquivalencia{}, the conventional threshold of a small
correlation~\cite{cohen1988power}.

\paragraph{Robustness}
Files from the same repository are not independent. We recompute the confidence interval
of RQ1 with a cluster bootstrap that resamples whole repositories (\FreplicasBootstrap{}
replicates), and repeat RQ1 without the files whose tests are affected by the nested class
defect.

\paragraph{Exploratory regime}
We ran \FtestesTotal{} statistical tests in total, all logged in the replication package,
including those that led nowhere. We present the study as exploratory: results that do not
survive the false discovery rate correction are reported as hypotheses, not findings.

\subsection{Reproducibility}
\label{sec:reproducibility}

The replication package contains the mirrors' mapping, the raw JNose reports, the
normalized tables, the audit spreadsheets, and one script per analysis. Every number in
this paper is a macro generated from the versioned data by a single script, and a checker
fails if any number in the documentation lacks a command that produces it.
